\FloatBarrier
\section{Endogenous design: additional results}\label{app:escDetails}\setcounter{equation}{0}
This appendix completes section \ref{sec:esc}: the three specifications of the cost of appendix \ref{app:esc:cost} across economies and across counterfactuals at $\rho = 1$, and the quadratic cost across intertemporal elasticities.

\subsection{Calibration across variants and countries}\label{app:escDetails:calibration}
Table \ref{table:US_ESC:variantsCountry} sets the three variants against the UK, France and the UK regrouped at the U.S.\ income percentiles, and adds the linear-quadratic cost with its shape pinned by France's corner instead of the UK's design. The own parameters say what each design asks of the cost. The UK's is within $16\%$ of the U.S.\ value under the quadratic cost, $7.264$ against $8.643$, and $41\%$ above it under the linear one, $2.414$ against $1.710$. France's is a threshold, the smallest value at which its electorate chooses the corner: no finite value under the quadratic cost, since the first unit of redistribution is free at the margin; $1.312$ under the linear cost, below the U.S.\ value; and $\lambda \geq 9.793$ under the linear-quadratic cost at the U.S.\ shape, nearly twice the U.S.\ value. Pinning the shape by France's corner instead takes $\kappa = 0.629$, nearly five times the UK's, and a linear parameter $\lambda_1 = 1.301$, about twice the UK's; the UK then chooses $0.397$ and would re-elect its design only at parameters about a fifth above the U.S.\ ones. The two pinnings are two points on one curve. Along the pairs $(\lambda,\kappa)$ that keep the U.S.\ on its design the marginal cost of redistribution at the U.S.\ design is held to first order, and the marginal cost at a design $\theta$ relative to it, $[\kappa M + \tilde{V}(1-\theta)]/[\kappa M + \tilde{V}(1-\theta^{\ast})]$ at the U.S.\ dispersions, falls with $\kappa$ for $\theta < \theta^{\ast}$ and rises for $\theta > \theta^{\ast}$; the UK, with the U.S.\ ratio $M/\tilde{V}$, therefore chooses a design that falls with the shape parameter, and France, with a larger ratio, one that rises, both monotonically, so no pair on the curve puts both on their designs. Of the three economies, the weight of the linear term matters for France alone. The U.S.\ and the UK have nearly the same $M/\tilde{V}$, $1.32$ and $1.31$, so with $\tilde{V}$ in place of $M$ on the linear term they bear nearly the same cost. France, at $1.93$, then chooses $0.749$, as under the quadratic cost, rather than $0.806$.

%% GENERATED by python/paper/build.py from results/ -- do not edit by hand.
%% Rebuild:  .venv\Scripts\python.exe python\paper\build.py --only US_ESC_VariantsCountry
%% Source:   results/esc/escCalibration{,_kappaFR}.csv, escCountry{,_kappaFR}.csv
\begin{table}[!htb]
\centering
\begin{threeparttable}
\caption{The cost specifications across economies}
\label{table:US_ESC:variantsCountry}
\renewcommand{\arraystretch}{1.25}
\footnotesize
\begin{tabularx}{\textwidth}{@{}>{\raggedright\arraybackslash}p{3.1cm}*{7}{Y}@{}}
\toprule
\textbf{Economy} & $\theta^{\ast}$ & \textbf{Tax rate} & $\tilde V$ & $M$ & $\theta$ \textbf{chosen} & \textbf{Own} $\lambda_1$ & \textbf{Own} $\lambda_2$ \\
\midrule 
\multicolumn{8}{l}{\textit{Linear-quadratic cost, $\kappa = 0.129$ pinned by the UK}} \\
U.S. & 0.738 & 14.4\% & 0.436 & 0.575 & 0.738 & 0.674 & 5.237 \\ % row: usa@sizeLQ
UK & 0.560 & 11.9\% & 0.223 & 0.292 & 0.560 & 0.674 & 5.237 \\ % row: gbr@sizeLQ
France & 1.000 & 21.3\% & 0.218 & 0.422 & 0.806 & $\geq 1.260$ & $\geq 9.793$ \\ % row: fra@sizeLQ
UK at U.S.\ income groups & 0.543 & 11.9\% & 0.065 & 0.239 & 0.846 & 0.411 & 3.196 \\[.5em] % row: ukus@sizeLQ
\multicolumn{8}{l}{\textit{Linear cost}} \\
U.S. & 0.738 & 14.4\% & 0.436 & 0.575 & 0.738 & 1.710 & 0 \\ % row: usa@sizeL
UK & 0.560 & 11.9\% & 0.223 & 0.292 & 0.000 & 2.414 & 0 \\ % row: gbr@sizeL
France & 1.000 & 21.3\% & 0.218 & 0.422 & 1.000 & $\geq 1.312$ & 0 \\ % row: fra@sizeL
UK at U.S.\ income groups & 0.543 & 11.9\% & 0.065 & 0.239 & 1.000 & 0.807 & 0 \\[.5em] % row: ukus@sizeL
\multicolumn{8}{l}{\textit{Quadratic cost}} \\
U.S. & 0.738 & 14.4\% & 0.436 & 0.575 & 0.738 & 0 & 8.643 \\ % row: usa@size
UK & 0.560 & 11.9\% & 0.223 & 0.292 & 0.625 & 0 & 7.264 \\ % row: gbr@size
France & 1.000 & 21.3\% & 0.218 & 0.422 & 0.749 & -- & -- \\ % row: fra@size
UK at U.S.\ income groups & 0.543 & 11.9\% & 0.065 & 0.239 & 0.644 & 0 & 6.519 \\[.5em] % row: ukus@size
\multicolumn{8}{l}{\textit{Linear-quadratic cost, $\kappa = 0.629$ pinned by France}} \\
U.S. & 0.738 & 14.4\% & 0.436 & 0.575 & 0.738 & 1.301 & 2.068 \\ % row: usa@sizeLQ_kappaFR
UK & 0.560 & 11.9\% & 0.223 & 0.292 & 0.397 & 1.580 & 2.512 \\ % row: gbr@sizeLQ_kappaFR
France & 1.000 & 21.3\% & 0.218 & 0.422 & 1.000 & $\geq 1.302$ & $\geq 2.070$ \\ % row: fra@sizeLQ_kappaFR
UK at U.S.\ income groups & 0.543 & 11.9\% & 0.065 & 0.239 & 1.000 & 0.674 & 1.072 \\ % row: ukus@sizeLQ_kappaFR
\bottomrule
\end{tabularx}
\begin{tablenotes}[flushleft]
\footnotesize
\item[] \textit{Note:} Log preferences, $\rho = 1$. One block per specification of the cost of table \ref{table:US_ESC:variants}, each with its own U.S.\ calibration; the last block pins the linear-quadratic cost\textquotesingle s $\kappa$ by France\textquotesingle s corner instead of the UK\textquotesingle s design. $\theta^{\ast}$ and the tax rate are each economy\textquotesingle s observed design and 2020 tax rate, those of table \ref{table:US:Calib} except the design of the UK at U.S.\ income groups, which is this table\textquotesingle s own, read off the UK workbook regrouped at the U.S.\ income percentiles; $\tilde V = \sum_i\gamma_i(y_i-1)^2/y_i$ and $M = \sum_i\gamma_i|y_i-1|$ are the dispersions of relative incomes that weight the quadratic and the linear term of the cost. ``$\theta$ chosen'' is the design in 2020 on the economy\textquotesingle s own freely simulated path under the block\textquotesingle s U.S.\ parameters, which the U.S.\ re-elects by construction; ``own'' are the parameters at which that path re-elects the observed design, with $\omega$ recalibrated at each trial value and, under the linear-quadratic cost, $\kappa = \lambda_1/\lambda_2$ held at the block\textquotesingle s value. France\textquotesingle s observed design is the corner $\theta = 1$, so its own values are the smallest at which its electorate chooses the corner ($\geq$); under the quadratic cost no finite value does (--). Each economy is its own calibration at $\rho = 1$ ($\beta$ imposed from the U.S.). With $\tilde V$ in place of $M$ on the linear term, at $\kappa = 0.170$ and with the U.S.\ design re-elected ($\lambda_1 = 0.889$, $\lambda_2 = 5.237$), France chooses $0.749$ under the U.S.\ parameters.
\end{tablenotes}
\end{threeparttable}
\end{table}


Regrouping the UK at the U.S.\ percentiles re-reads its observed design as $0.543$ and compresses its relative incomes: the dispersion $\tilde{V}$ falls from $0.223$ to $0.065$, the amount carried $M$ only from $0.292$ to $0.239$.\footnote{The regrouping is a linear fit of the UK's group means on the cumulative population shares rather than a recount of the survey (appendix \ref{app:US}), and the fit compresses the distribution. Since the UK's own low and medium groups, $29\%$ and $58\%$ of households, have mean incomes of half the mean and the mean, the bottom $58\%$ average at most $0.75$ of mean income and the top $24\%$ at least $1.6$, so an exact regrouping would have $\tilde{V}$ of at least $0.10$ and a lower $M/\tilde{V}$, the ratio through which a linear term moves the regrouped UK's design.} Under the quadratic cost the choice moves little, to $0.644$ against $0.625$ at the UK's own cuts, since the redistributive stakes and the cost shrink together. Under a cost with a linear term it moves a long way, to $0.846$ against $0.560$ under the linear-quadratic cost and to the Bismarckian corner under the linear cost and under the shape pinned by France: its $M/\tilde{V}$ is $3.7$ against $1.31$ at its own cuts, so the linear term weighs far more on it relative to its redistributive stakes, the mechanism by which France's income distribution moves the U.S.\ design (section \ref{sec:esc:results}).

\FloatBarrier
\subsection{Counterfactuals across variants}\label{app:escDetails:counterfactuals}
Table \ref{table:US_ESC:variantsShocks} reads every scenario under the three variants, with the savings rate and the workweek. Mild ageing moves the chosen design a little more than half as far as acute ageing under the linear-quadratic and the quadratic cost, to $0.811$ and $0.784$, and to the Bismarckian corner under the linear cost. Under the linear-quadratic cost France's income distribution moves the design further than mild ageing does and nearly as far as acute ageing, $0.850$ against $0.874$, whereas under the quadratic cost it moves it less than mild ageing, $0.772$ against $0.784$. Choosing the design adds to the rise in the tax rate under both ageing scenarios and every variant, by at most $0.4$ p.p.\ under the linear-quadratic cost and $1.3$ p.p.\ under the linear one.

%% GENERATED by python/paper/build.py from results/ -- do not edit by hand.
%% Rebuild:  .venv\Scripts\python.exe python\paper\build.py --only US_ESC_VariantsShocks
%% Source:   results/esc/escExperiments.csv
\begin{table}[!htb]
\centering
\begin{threeparttable}
\caption{Endogenous pension design across counterfactuals and cost specifications: all readings}
\label{table:US_ESC:variantsShocks}
\renewcommand{\arraystretch}{1.25}
\footnotesize
\begin{tabularx}{\textwidth}{@{}>{\raggedright\arraybackslash}X*{7}{c}@{}}
\toprule
 & & \multicolumn{2}{c}{\textbf{Tax rate}} & \multicolumn{2}{c}{\textbf{Savings rate}} & \multicolumn{2}{c}{\textbf{Workweek}} \\
\cmidrule(lr){3-4}\cmidrule(lr){5-6}\cmidrule(lr){7-8}
\textbf{Scenario} & $\theta$ \textbf{chosen} & \textbf{Pinned} & \textbf{Chosen} & \textbf{Pinned} & \textbf{Chosen} & \textbf{Pinned} & \textbf{Chosen} \\
\midrule 
\multicolumn{8}{l}{\textit{Linear-quadratic cost}} \\
Baseline & 0.738 & 14.43\% & 14.43\% & 15.54\% & 15.41\% & 39.39 & 39.51 \\ % row: baseline@sizeLQ
Mild ageing & 0.811 & 18.08\% & 18.22\% & $-0.77$ p.p. & $-0.97$ p.p. & 39.18 & 39.29 \\ % row: mild@sizeLQ
Acute ageing & 0.874 & 22.29\% & 22.73\% & $-1.76$ p.p. & $-2.05$ p.p. & 38.90 & 38.96 \\ % row: acute@sizeLQ
Income distribution & 0.850 & 13.81\% & 13.91\% & $+0.28$ p.p. & $+0.03$ p.p. & 40.23 & 40.45 \\ % row: frincome@sizeLQ
Voting patterns & 0.622 & 15.25\% & 15.25\% & $-0.15$ p.p. & $-0.12$ p.p. & 39.21 & 39.20 \\ % row: frvoting@sizeLQ
Income distribution and voting patterns & 0.641 & 14.41\% & 14.41\% & $+0.06$ p.p. & $+0.04$ p.p. & 40.10 & 40.13 \\ % row: frboth@sizeLQ
All French characteristics & 0.641 & 14.41\% & 14.41\% & $+0.06$ p.p. & $+0.04$ p.p. & 34.95 & 34.98 \\ % row: frall@sizeLQ
France & 1.000 & 21.29\% & -- & $-1.97$ p.p. & -- & 35.24 & -- \\[.5em] % row: france@sizeLQ
\multicolumn{8}{l}{\textit{Linear cost}} \\
Baseline & 0.738 & 14.43\% & 14.43\% & 15.61\% & 15.16\% & 39.39 & 39.76 \\ % row: baseline@sizeL
Mild ageing & 1.000 & 17.93\% & 18.57\% & $-0.41$ p.p. & $-1.01$ p.p. & 39.22 & 39.41 \\ % row: mild@sizeL
Acute ageing & 1.000 & 21.97\% & 23.24\% & $-1.36$ p.p. & $-2.09$ p.p. & 38.99 & 39.02 \\ % row: acute@sizeL
Income distribution & 1.000 & 13.85\% & 14.51\% & $+0.58$ p.p. & $-0.01$ p.p. & 40.22 & 40.42 \\ % row: frincome@sizeL
Voting patterns & 0.047 & 15.22\% & 15.01\% & $+0.18$ p.p. & $+0.44$ p.p. & 39.22 & 39.15 \\ % row: frvoting@sizeL
Income distribution and voting patterns & 0.743 & 14.44\% & 14.44\% & $+0.37$ p.p. & 0.00 p.p. & 40.09 & 40.45 \\ % row: frboth@sizeL
All French characteristics & 0.743 & 14.44\% & 14.44\% & $+0.37$ p.p. & 0.00 p.p. & 34.93 & 35.24 \\ % row: frall@sizeL
France & 1.000 & 21.29\% & -- & $-1.72$ p.p. & -- & 35.24 & -- \\[.5em] % row: france@sizeL
\multicolumn{8}{l}{\textit{Quadratic cost}} \\
Baseline & 0.738 & 14.43\% & 14.43\% & 15.50\% & 15.41\% & 39.39 & 39.47 \\ % row: baseline@size
Mild ageing & 0.784 & 18.18\% & 18.27\% & $-0.84$ p.p. & $-0.96$ p.p. & 39.15 & 39.22 \\ % row: mild@size
Acute ageing & 0.825 & 22.51\% & 22.80\% & $-1.85$ p.p. & $-2.04$ p.p. & 38.83 & 38.87 \\ % row: acute@size
Income distribution & 0.772 & 13.78\% & 13.80\% & $+0.25$ p.p. & $+0.16$ p.p. & 40.24 & 40.34 \\ % row: frincome@size
Voting patterns & 0.653 & 15.27\% & 15.26\% & $-0.20$ p.p. & $-0.18$ p.p. & 39.21 & 39.20 \\ % row: frvoting@size
Income distribution and voting patterns & 0.627 & 14.40\% & 14.44\% & $+0.03$ p.p. & $+0.05$ p.p. & 40.10 & 40.06 \\ % row: frboth@size
All French characteristics & 0.627 & 14.40\% & 14.44\% & $+0.03$ p.p. & $+0.05$ p.p. & 34.96 & 34.93 \\ % row: frall@size
France & 1.000 & 21.29\% & -- & $-1.97$ p.p. & -- & 35.24 & -- \\ % row: france@size
\bottomrule
\end{tabularx}
\begin{tablenotes}[flushleft]
\footnotesize
\item[] \textit{Note:} Log preferences, $\rho = 1$. One block per specification of the cost, each at its U.S.\ calibration (table \ref{table:US_ESC:variants}). Every row is a separate equilibrium path read in 2020, with the changed characteristic holding over the whole horizon and the political choice binding from the first period. ``Pinned'' holds the design at $\theta^{\ast}$ and ``chosen'' lets the electorate choose it one period in advance, both under the cost. The ageing scenarios are those of section \ref{sec:oecd}; the two composite rows impose France\textquotesingle s income distribution and voting patterns together and, in the second, also its level of $X$. The France row is France\textquotesingle s own calibrated path at its observed design $\theta = 1$, where none of the three costs bites, and has no chosen reading. The savings rate is savings relative to GDP: its level on the baseline row and, on every other row, the change against the baseline with the design chosen, in percentage points. The workweek is in hours.
\end{tablenotes}
\end{threeparttable}
\end{table}


The savings rate and the workweek of 2020 respond to the design chosen in 2020 for the retirement of the young: a more earnings-related benefit lowers the marginal wedge of the flat component on their hours and, with less of the revenue lost, raises what they expect from the pension, so they work more and save less. On the baseline that design is above the observed one under every variant (figure \ref{fig:US_ESC:path}), so even there the savings rate is lower with the design chosen than with it pinned and the workweek longer, by $0.1$ p.p.\ of GDP and $0.1$ hours under the linear-quadratic cost and by $0.5$ p.p.\ of GDP and $0.4$ hours under the linear cost. Ageing widens the gap in the savings rate under every variant, to $0.3$ p.p.\ under the linear-quadratic cost and $0.7$ p.p.\ under the linear one at acute ageing, and narrows the gap in the workweek, since the chosen design then comes with a higher tax rate. France's voting patterns, under which the electorate chooses flatter designs, reverse both gaps.

Under the linear cost France's income distribution and voting patterns, which alone send the design to a corner or next to it, together leave it at $0.743$, beside the U.S.\ design; under the other two variants the voting patterns prevail (section \ref{sec:esc:results}). Adding France's taste for leisure shortens the workweek by about five hours under every variant and leaves the design and every rate as they are, as it does with the design given (section \ref{sec:oecd}). With all of France's household characteristics the U.S.\ keeps a tax rate of $14.4\%$, against $21.3\%$ on France's own path, and chooses $0.641$ under the linear-quadratic cost, against France's $0.806$ under the same parameters (table \ref{table:US_ESC:variantsCountry}): with a political weight of the old no higher than the U.S.\ one, what sets France's system apart is its demography rather than its households. On the UK as host, at its own parameters, France's voting patterns move the design toward flat benefits under every variant and prevail over its income distribution, which leaves the UK's design within $0.01$ of its observed value except under the linear cost (\oa{esc-variants}).

\FloatBarrier
\subsection{The role of the intertemporal elasticity}\label{app:escDetails:rho}
Tables \ref{table:US_ESC:calibration} and \ref{table:US_ESC:summary} hold the quadratic cost at $\rho = 0.5$, $1$ and $2$. As its parameter falls with $\rho$, so does the share of revenue lost at the observed design, from $3.9\%$ at $\rho = 0.5$ to $0.4\%$ at $\rho = 2$, and at a flat benefit, from $43.7\%$ to $5.3\%$. A higher elasticity strengthens the forward-looking channels that hold the choice away from the Beveridgean corner, so less cost is needed to keep the U.S.\ electorate at its observed design. Without any cost the same channels show in the timing of the choice: chosen with the tax, the design is the Beveridgean corner at every elasticity checked, while chosen once and for all, which enters every channel at once, it is the Beveridgean corner up to $\rho = 1.3$ and the Bismarckian corner from $\rho = 1.4$ on (\oa{esc-timing}).

%% GENERATED by python/paper/build.py from results/ -- do not edit by hand.
%% Rebuild:  .venv\Scripts\python.exe python\paper\build.py --only US_ESC_Calibration
%% Source:   results/esc/escCalibration{,CRRA}.csv
\begin{table}[!htb]
\centering
\begin{threeparttable}
\caption{The calibrated cost of redistribution for the U.S.}
\label{table:US_ESC:calibration}
\renewcommand{\arraystretch}{1.25}
\begin{tabularx}{.8\textwidth}{YYYYY}
\toprule
$\rho$ & $\lambda$ & $\omega$ & $1 - f(\theta^{\ast})$ & $1 - f(0)$ \\
\midrule 
0.5 & 18.267 & 2.50 & 3.9\% & 43.7\% \\ % row: rho@0.5
1.0 & 8.643 & 1.51 & 1.8\% & 23.8\% \\ % row: rho@1.0
2.0 & 1.724 & 1.20 & 0.4\% & 5.3\% \\ % row: rho@2.0
\bottomrule
\end{tabularx}
\begin{tablenotes}[flushleft]
\footnotesize
\item[] \textit{Note:} The cost is that of equation \eqref{eq:esc:budget}, with $\theta^{\ast} = 0.738$ and $\tilde V = 0.436$ at every $\rho$. $1 - f(\theta^{\ast})$ and $1 - f(0)$ are the shares of revenue lost in 2020 at the observed design and at a flat benefit, both at the 2020 tax rate. $\beta$ and $\omega$ are recalibrated with $\lambda$.
\end{tablenotes}
\end{threeparttable}
\end{table}


Acute ageing moves the chosen design about equally at every $\rho$ (section \ref{sec:esc:results}), and choosing the design adds $0.1$ to $0.3$ p.p.\ to the rise in the tax rate under it. France's income distribution moves the design toward flat benefits at $\rho = 0.5$, to $0.707$, and toward earnings-related ones at $\rho = 1$ and $2$; its voting patterns move it toward flat benefits at every $\rho$, furthest at $\rho = 2$, and the composite of the two follows them, to $0.667$, $0.627$ and $0.262$. At $\rho = 2$, where the cost is thinnest, choosing the design also raises the tax rate, by $0.7$ p.p.\ under France's voting patterns and $0.4$ p.p.\ under the composite, while at $\rho \leq 1$ it moves the tax rate by at most $0.1$ p.p.\ in these scenarios. Across $\rho$ under the quadratic cost, as across the variants at $\rho = 1$, ageing makes the chosen design more Bismarckian everywhere, while the response to inequality changes direction with $\rho$ and size with the variant.

%% GENERATED by python/paper/build.py from results/ -- do not edit by hand.
%% Rebuild:  .venv\Scripts\python.exe python\paper\build.py --only US_ESC_Summary
%% Source:   results/esc/escExperiments.csv
\begin{table}[!htb]
\centering
\begin{threeparttable}
\caption{Endogenous pension design across counterfactuals and intertemporal elasticities in the U.S.}
\label{table:US_ESC:summary}
\renewcommand{\arraystretch}{1.25}
\footnotesize
\begin{tabularx}{\textwidth}{@{}>{\raggedright\arraybackslash}p{2.4cm}*{9}{@{\hspace{9pt}}Y}@{}}
\toprule
 & \multicolumn{3}{c}{$\rho = 0.5$} & \multicolumn{3}{c}{$\rho = 1.0$} & \multicolumn{3}{c}{$\rho = 2.0$} \\
\cmidrule(lr){2-4}\cmidrule(lr){5-7}\cmidrule(lr){8-10}
\textbf{Scenario} & $\theta$ \textbf{chosen} & \textbf{Tax, pinned} & \textbf{Tax, chosen} & $\theta$ \textbf{chosen} & \textbf{Tax, pinned} & \textbf{Tax, chosen} & $\theta$ \textbf{chosen} & \textbf{Tax, pinned} & \textbf{Tax, chosen} \\
\midrule 
Baseline & 0.738 & 14.43\% & 14.41\% & 0.738 & 14.43\% & 14.43\% & 0.738 & 14.43\% & 14.44\% \\ % row: baseline
Acute ageing & 0.816 & 23.01\% & 23.32\% & 0.825 & 22.51\% & 22.80\% & 0.826 & 21.79\% & 21.92\% \\ % row: acute
Income distribution & 0.707 & 13.71\% & 13.67\% & 0.772 & 13.78\% & 13.80\% & 0.918 & 13.85\% & 14.09\% \\ % row: frincome
Voting patterns & 0.687 & 14.73\% & 14.62\% & 0.653 & 15.27\% & 15.26\% & 0.335 & 16.08\% & 16.73\% \\ % row: frvoting
French characteristics & 0.667 & 13.98\% & 13.96\% & 0.627 & 14.40\% & 14.44\% & 0.262 & 15.04\% & 15.45\% \\ % row: frboth
\bottomrule
\end{tabularx}
\begin{tablenotes}[flushleft]
\footnotesize
\item[] \textit{Note:} The deadweight cost $f(\theta_t, \tau_t)$ of equation \eqref{eq:esc:budget} is the quadratic cost, $f = \exp\{-\tfrac12 \lambda \tau_t \tilde V (1-\theta_t)^2\}$: quadratic in the implicit tax the flat component levies and scaled by the size of the system, with $\lambda$ calibrated per $\rho$ (table \ref{table:US_ESC:calibration}). Every row is a separate equilibrium path read in 2020, with the changed characteristic holding over the whole horizon and the political choice binding from the first period. ``Pinned'' holds the design at the U.S.\ value $\theta^{\ast}$ and ``chosen'' lets the electorate choose it one period in advance, both under the cost. The composite imposes France\textquotesingle s income distribution and voting patterns together. France\textquotesingle s level of $X$, which all French characteristics add, leaves every cell of the row unchanged except the last digit of the tax at $\rho = 0.5$, so up to that digit it is also the design and the tax under all French characteristics. The savings and workweek readings are in the \oa{esc-us}.
\end{tablenotes}
\end{threeparttable}
\end{table}


The UK's own quadratic cost falls less steeply with $\rho$ than the U.S.\ one: $16$ to $17\%$ below it at $\rho \leq 1$, it is $61\%$ above it at $\rho = 2$ (\oa{esc-uk}). The Frisch elasticity, through which the design acts on hours, matters little: recalibrated at $\xi = 0.2$ and $0.4$, the cost parameter moves by about one percent and the chosen designs on the baseline path and under acute ageing by at most $0.001$, while the tax rate under acute ageing moves by about $0.4$ p.p., lower at the higher elasticity (\oa{esc-path}). The linear and the linear-quadratic costs are solved under log preferences only, so the comparison across $\rho$ is the quadratic cost's.
